\chapter{Measuring success: ROC-AUC}
\label{ch:auc}

\begin{tldr}
\begin{itemize}
\item \textbf{Accuracy is useless} when positives are rare: ``always no'' scores 98\,\% on MCL.
\item \textbf{AUC} = the share of (sick, healthy) pairs that the model puts in the right order. 0.5 = random, 1 = perfect.
\item Only the \textbf{order} of the predictions matters, never their exact values.
\item With few test cases AUC is \textbf{noisy}: on our 58 gold studies, $\pm$0.02 macro is luck.
\end{itemize}
\end{tldr}

\section{Why not just count correct answers?}

\begin{example}
In our report labels, only about 2\,\% of studies have a high-grade MCL tear. A ``model'' that always answers
\emph{no} is right 98\,\% of the time. Its accuracy is excellent, and it is completely useless: it finds no tears.
\end{example}

We need a score that rewards \emph{finding the sick ones} and does not care how rare they are. That is ROC-AUC.

\section{AUC as a ranking game}

Give the model 8 studies. It returns a score for each; sort them from highest to lowest:

\begin{figure}[H]
\centering
\begin{tikzpicture}[x=1mm,y=1mm]
  \draw[muted, -{Stealth}] (0,0) -- (132,0) node[right, font=\scriptsize] {score};
  \foreach \s/\p/\n in {95/1/S1, 85/1/S2, 70/0/S3, 60/1/S4, 40/0/S5, 30/0/S6, 20/0/S7, 10/0/S8} {
    \pgfmathsetmacro{\x}{\s*1.25}
    \ifnum\p=1
      \fill[them] (\x,6) circle (2.6);
      \node[font=\scriptsize\bfseries, text=white] at (\x,6) {+};
    \else
      \fill[us] (\x,6) circle (2.6);
      \node[font=\scriptsize\bfseries, text=white] at (\x,6) {--};
    \fi
    \node[font=\scriptsize] at (\x,11.5) {\n};
    \node[font=\scriptsize, text=muted] at (\x,-3.5) {0.\s};
  }
  % the one wrong pair
  \draw[warnline, line width=0.8pt, -{Stealth}] (75,15.5) to[bend left=35] (87.5,15.5);
  \node[font=\scriptsize, text=warnline, align=center] at (81,22) {the only pair in the wrong order:\\ S3 (healthy) above S4 (sick)};
\end{tikzpicture}
\caption{Eight studies on the model's score axis (higher score to the right). Orange~+ = has the finding,
blue~-- = does not.}
\label{fig:rankgame}
\end{figure}

Now form every pair of one sick and one healthy study: $3 \times 5 = 15$ pairs. Count the pairs where the sick one
has the higher score:
\begin{itemize}
\item S1 beats all 5 healthy ones: 5 \quad S2 beats all 5: 5 \quad S4 beats S5--S8 but not S3: 4.
\end{itemize}
\[
\text{AUC} = \frac{5+5+4}{15} = \frac{14}{15} = 0.933
\]

\begin{keyidea}
\textbf{AUC is the probability that a random sick study gets a higher score than a random healthy one.}
It does not depend on how rare the finding is, and it does not depend on the scores themselves, only on
their order.
\end{keyidea}

\section{The ROC curve}

The same number has a picture. Move a threshold from the top score down. At each step, count the sick studies above
the threshold (\term{true-positive rate}, TPR) and the healthy ones above it (\term{false-positive rate}, FPR):

\begin{figure}[H]
\centering
\begin{tikzpicture}
\begin{axis}[width=6.3cm, height=6.3cm, xmin=0, xmax=1, ymin=0, ymax=1.02, xlabel={false-positive rate},
    ylabel={true-positive rate}, xtick={0,0.2,...,1}, ytick={0,0.333,0.667,1}, yticklabels={0,1/3,2/3,1},
    xlabel style={font=\footnotesize}, ylabel style={font=\footnotesize}, grid=major]
  \addplot[fill=us!15, draw=none] coordinates {(0,0) (0,0.667) (0.2,0.667) (0.2,1) (1,1) (1,0)} \closedcycle;
  \addplot[us, line width=1.5pt, mark=*, mark size=1.6pt] coordinates
     {(0,0) (0,0.333) (0,0.667) (0.2,0.667) (0.2,1) (0.4,1) (0.6,1) (0.8,1) (1,1)};
  \addplot[muted, dashed] coordinates {(0,0) (1,1)};
  \node[font=\footnotesize] at (axis cs:0.6,0.45) {area = 0.933};
  \node[font=\scriptsize, anchor=west] at (axis cs:0.02,0.28) {S1};
  \node[font=\scriptsize, anchor=west] at (axis cs:0.02,0.62) {S2};
  \node[font=\scriptsize, anchor=south] at (axis cs:0.1,0.67) {S3};
  \node[font=\scriptsize, anchor=west] at (axis cs:0.21,0.9) {S4};
\end{axis}
\end{tikzpicture}\hspace{6mm}
\begin{minipage}[b]{0.42\linewidth}\small
Start at (0,0). Walk down the sorted list:
\begin{itemize}
\item a \textcolor{them}{\textbf{sick}} study: step \textbf{up} by 1/3;
\item a \textcolor{us}{\textbf{healthy}} study: step \textbf{right} by 1/5.
\end{itemize}
The area under the staircase is the AUC. A perfect model goes straight up, then right (area 1). A random model
follows the dashed diagonal (area 0.5).\par\bigskip
\end{minipage}
\caption{The ROC curve of the 8-study example. ROC = ``receiver operating characteristic'', a name from radar
engineering; AUC = ``area under the curve''.}
\end{figure}

\section{Real curves from our model}

These are the actual ROC curves of our v4 model (5 folds, EfficientNet-B0) on the 58 gold studies:

\begin{figure}[H]
\centering
\begin{tikzpicture}
\begin{axis}[width=7.2cm, height=7.2cm, xmin=0, xmax=1, ymin=0, ymax=1, xlabel={false-positive rate},
    ylabel={true-positive rate}, xlabel style={font=\footnotesize}, ylabel style={font=\footnotesize}, grid=major,
    legend style={at={(0.97,0.03)}, anchor=south east}]
  \addplot[them, line width=1.4pt, const plot] table[x=fpr, y=tpr] {figdata/roc_eff.dat};
  \addplot[us, line width=1.4pt, const plot] table[x=fpr, y=tpr] {figdata/roc_acl.dat};
  \addplot[third, line width=1.4pt, const plot] table[x=fpr, y=tpr] {figdata/roc_syn.dat};
  \addplot[muted, dashed] coordinates {(0,0) (1,1)};
  \legend{{Effusion: \auceff}, {ACL: \aucacl}, {Synovitis: \aucsyn}}
\end{axis}
\end{tikzpicture}
\caption{Effusion and ACL are found well; Synovitis is our weakest finding. With only \goldPosSyn{} synovitis
positives among the 58 studies the staircase has large steps: every single study moves the curve.}
\label{fig:realroc}
\end{figure}

\section{Twelve AUCs become one: the macro average}

\begin{figure}[H]
\centering
\begin{tikzpicture}
\begin{axis}[xbar, width=10cm, height=7cm, bar width=2.6mm, xmin=0.6, xmax=1.02,
    symbolic y coords={Fracture,Contusion,Baker's,Synovitis,Effusion,PF OA,Lat.~OA,Med.~OA,Lat.~meniscus,Med.~meniscus,MCL,ACL},
    ytick=data, y tick label style={font=\footnotesize}, xtick={0.6,0.7,0.8,0.9,1.0},
    xmajorgrids, nodes near coords, nodes near coords style={font=\scriptsize, anchor=west},
    point meta=explicit symbolic, enlarge y limits=0.05]
  \addplot[barus] coordinates {(0.934,ACL) [0.934] (0.887,MCL) [0.887] (0.900,Med.~meniscus) [0.900]
     (0.784,Lat.~meniscus) [0.784] (0.977,Med.~OA) [0.977] (0.820,Lat.~OA) [0.820] (0.900,PF OA) [0.900]
     (0.955,Effusion) [0.955] (0.711,Synovitis) [0.711] (0.980,Baker's) [0.980] (0.901,Contusion) [0.901]
     (0.865,Fracture) [0.865]};
  \addplot[them, line width=1.2pt, dashed, nodes near coords={}] coordinates {(0.884,Fracture) [] (0.884,ACL) []};
\end{axis}
\end{tikzpicture}
\caption{Our v4 model on the 58 gold studies, per finding. The dashed line is the macro average, \goldMacro.
Synovitis, Lateral meniscus and Lateral OA pull it down most.}
\label{fig:perfinding}
\end{figure}

\begin{why}
The competition uses the \emph{macro} average (each finding counts 1/12) so that a model cannot ignore the rare
findings. Improving Synovitis from 0.71 to 0.81 lifts the score by $0.10/12 \approx 0.008$, as much as improving every
finding by 0.008.
\end{why}

\section{Only the order matters}

Because AUC only looks at the order of scores:
\begin{itemize}
\item multiplying all predictions by 10, squaring them or passing them through any increasing function changes
      nothing;
\item a model can be badly \emph{calibrated} (its 0.8 does not mean ``80\,\% chance'') and still have a perfect AUC;
\item to combine models, you can average their \textbf{ranks} instead of their probabilities (Chapter~\ref{ch:submission}).
\end{itemize}

\begin{watchout}
Calibration does not matter for the \emph{score}, but it matters for \emph{training labels}. A label of 0.98 that
should be 0.6 teaches the model to be confidently wrong. The winners judged their labels on calibration too
(Chapter~\ref{ch:winlabels}).
\end{watchout}

\section{How noisy is an AUC?}

\begin{figure}[H]
\centering
\begin{tikzpicture}
\begin{axis}[ybar interval, width=10cm, height=4.8cm, xmin=0.82, xmax=0.94, ymin=0, ymax=330, xlabel={macro AUC on a resampled gold set},
    ytick=\empty, axis y line=none, axis x line*=bottom, xlabel style={font=\footnotesize},
    xtick={0.82,0.84,...,0.94}, xticklabel style={/pgf/number format/fixed, /pgf/number format/precision=2}]
  \addplot[fill=us!70, draw=white] table[x=left, y=n] {figdata/gold_bootstrap.dat};
  \draw[them, line width=1.4pt] (axis cs:0.884,0) -- (axis cs:0.884,300);
  \node[font=\scriptsize, text=them!85!black, anchor=west] at (axis cs:0.886,290) {measured: \goldMacro};
\end{axis}
\end{tikzpicture}
\caption{The same predictions, scored on 2,000 random resamples of the 58 gold studies (a \term{bootstrap}). The
middle 90\,\% of the results spans \bootLo--\bootHi. A difference of 0.01 between two models on gold tells you almost
nothing.}
\label{fig:bootstrap}
\end{figure}

\begin{innumbers}
\begin{itemize}
\item 58 gold studies: macro AUC uncertain by about $\pm$0.02; a single finding by $\pm$0.05 to $\pm$0.09 (MCL has only
      9 positives).
\item Public leaderboard (\ca 390 studies): about $\pm$0.005.
\item Out-of-fold predictions on 4,349 report-labelled studies: much more stable, but measured against
      \emph{our own} labels, which have their own errors.
\end{itemize}
\end{innumbers}

\begin{keyidea}
Every score is measured against \emph{something}: the 58 gold labels (trustworthy but few), report labels (many
but noisy) or the public leaderboard (hidden, limited submissions). Use them together and be suspicious of small
differences. The winners spent real effort choosing their ruler (Chapter~\ref{ch:winmodels}).
\end{keyidea}

\begin{checkpoint}
\begin{itemize}
\done I can compute an AUC by counting pairs.
\done I can draw a ROC curve from a sorted list.
\done I know why a 0.01 gain on 58 gold studies is not evidence.
\quiz Q6.1: A model outputs 0.51 for every sick study and 0.49 for every healthy one. What is its AUC?
\quiz Q6.2: You add 0.3 to every prediction of one finding. What happens to the AUC?
\quiz Q6.3: Which matters more for the final score: raising Baker's from 0.98 to 0.99, or Synovitis from 0.71 to
      0.75?
\end{itemize}
\end{checkpoint}
